\documentclass[a4paper,12pt]{article}
\usepackage[utf8]{inputenc}
\usepackage[T1]{fontenc}
\usepackage{lmodern}
\usepackage{amsmath}
\usepackage{geometry}
\usepackage{booktabs}
\usepackage{graphicx}
\usepackage{enumitem}
\usepackage{xcolor}
\geometry{margin=2.5cm}
\setlength{\parindent}{0pt}
\setlength{\parskip}{0.6em}

\title{IoT - Challenge 3 - Part 2}
\author{Raul Agolli - Team Leader \\\\ Gabriele Lorenzo Singe}
\date{May 2026}

\begin{document}

\maketitle

\section*{Exercise EQ1}

We consider a LoRaWAN network with $N=40$ nodes, each generating packets according to a Poisson process with intensity $\lambda = 2$ packets/min.

The total traffic offered to the channel is:
\[
\Lambda = N \lambda = 40 \cdot 2 = 80 \text{ packets/min} = \frac{80}{60} \approx 1.33 \text{ packets/s}
\]

According to the LoRaWAN Class~A specification, the uplink access is ALOHA-like. Hence, assuming an asynchronous ALOHA model, the probability of successful transmission is:
\[
P_{\text{succ}} = e^{-2\Lambda T_{\text{air}}}
\]

We require:
\[
P_{\text{succ}} \ge 0.75
\]

Thus:
\[
e^{-2\Lambda T_{\text{air}}} \ge 0.75
\]

Taking the logarithm:
\[
T_{\text{air}} \le \frac{-\ln(0.75)}{2\Lambda}
\]

Substituting:
\[
T_{\text{air}} \le \frac{0.2877}{2 \cdot 1.33} \approx 0.108 \text{ s} = 107.9 \text{ ms}
\]

Using the TTN airtime calculator (EU868, BW = 125 kHz, payload = 20 bytes):

\[
\begin{array}{c|c}
\text{SF} & T_{\text{air}} \\
\hline
7 & 71.9\ \text{ms} \\
8 & 133.6\ \text{ms} \\
9 & 246.8\ \text{ms}
\end{array}
\]

Only SF7 satisfies the constraint:
\[
71.9\ \text{ms} \le 107.9\ \text{ms}
\]

\[
\boxed{\text{SF} = 7}
\]

\section*{Exercise EQ2}

The observed performance is significantly lower than expected and highly non-uniform across nodes. This indicates that the problem is not only due to collisions, but also to heterogeneous link conditions (e.g., distance, path loss, or SNR differences).

\begin{itemize}
    \item Decreasing the number of nodes would reduce collisions globally, but would not explain the strong performance asymmetry.
    \item Increasing the spreading factor would improve sensitivity, but also increase airtime, thus worsening collisions.
    \item Moving nodes closer to the gateway improves link quality (higher SNR, lower path loss), directly addressing the nodes with poor performance.
\end{itemize}

Therefore, the most appropriate action is:
\[
\boxed{\text{Move the nodes closer to the gateway}}
\]
\end{document}